% =====================================================================
% Abbildung: ROS2-Node-Graph des autonomen Indoor-Roboters
% Passend zu Kapitel 15, Abschnitt "Kommunikationsstruktur"
%
% Einbindung im Buch, z. B. in chapters/15_ros2_architektur.tex:
%
%   \begin{figure}[H]
%     \centering
%     \resizebox{\textwidth}{!}{% =====================================================================
% Abbildung: ROS2-Node-Graph des autonomen Indoor-Roboters
% Passend zu Kapitel 15, Abschnitt "Kommunikationsstruktur"
%
% Einbindung im Buch, z. B. in chapters/15_ros2_architektur.tex:
%
%   \begin{figure}[H]
%     \centering
%     \resizebox{\textwidth}{!}{% =====================================================================
% Abbildung: ROS2-Node-Graph des autonomen Indoor-Roboters
% Passend zu Kapitel 15, Abschnitt "Kommunikationsstruktur"
%
% Einbindung im Buch, z. B. in chapters/15_ros2_architektur.tex:
%
%   \begin{figure}[H]
%     \centering
%     \resizebox{\textwidth}{!}{% =====================================================================
% Abbildung: ROS2-Node-Graph des autonomen Indoor-Roboters
% Passend zu Kapitel 15, Abschnitt "Kommunikationsstruktur"
%
% Einbindung im Buch, z. B. in chapters/15_ros2_architektur.tex:
%
%   \begin{figure}[H]
%     \centering
%     \resizebox{\textwidth}{!}{\input{figures/fig_kap15_node_graph}}
%     \caption{Node-Graph des autonomen Indoor-Roboters mit den wichtigsten
%       Topics zwischen Sensor-, Verarbeitungs-, Navigations- und
%       Aktuatorik-Nodes.}
%     \label{fig:kap15-node-graph}
%   \end{figure}
%
% Benoetigte Pakete (im Buch bereits in main.tex geladen): tikz, graphicx
% Benoetigte TikZ-Bibliotheken: positioning, arrows.meta, shadows, calc
% =====================================================================

\input{figures/fig_kap15_node_graph.tikz}
}
%     \caption{Node-Graph des autonomen Indoor-Roboters mit den wichtigsten
%       Topics zwischen Sensor-, Verarbeitungs-, Navigations- und
%       Aktuatorik-Nodes.}
%     \label{fig:kap15-node-graph}
%   \end{figure}
%
% Benoetigte Pakete (im Buch bereits in main.tex geladen): tikz, graphicx
% Benoetigte TikZ-Bibliotheken: positioning, arrows.meta, shadows, calc
% =====================================================================

\begin{tikzpicture}[
    scale=0.74, transform shape,
    every node/.style={font=\small},
    node distance=10mm and 14mm,
    ros/.style={
        rectangle, rounded corners=2pt, draw=orange!60!black, thick,
        minimum width=32mm, minimum height=10mm,
        align=center, fill=orange!15, drop shadow
    },
    flow/.style={-{Stealth[length=2.2mm]}, thick, orange!70!black},
    lbl/.style={font=\footnotesize, fill=white, inner sep=1pt}
]

% --- Spalte 1: Sensor-Nodes (links) ---
\node[ros] (camera)  at (0, 6)  {\texttt{camera\_node}};
\node[ros] (lidar)   at (0, 3)  {\texttt{lidar\_node}};
\node[ros] (battery) at (0, 0)  {\texttt{battery\_monitor\_}\\\texttt{node}};

% --- Spalte 2: Verarbeitung ---
\node[ros] (detector) at (5.6, 6) {\texttt{object\_detector\_}\\\texttt{node}};
\node[ros] (slam)     at (5.6, 3) {\texttt{slam\_node}};

% --- Spalte 3: zentral ---
\node[ros] (nav) at (10.2, 3.5) {\texttt{navigation\_node}};

% --- Spalte 4: rechts ---
\node[ros] (motion) at (14.8, 5)   {\texttt{motion\_control\_}\\\texttt{node}};
\node[ros] (ui)      at (14.8, 0.5) {\texttt{user\_interface\_}\\\texttt{node}};

% --- zusätzlicher Node (Status) ---
\node[ros] (state) at (10.2, -1.5) {\texttt{robot\_state\_node}};

% --- Kanten ---
\draw[flow] (camera.east)  -- node[lbl, midway, above] {/camera/image\_raw} (detector.west);
\draw[flow] (lidar.east)   -- node[lbl, midway, above] {/scan}              (slam.west);
\draw[flow] (slam.east)    -- node[lbl, midway, above] {/map}               (nav.north west);
\draw[flow] (nav.east)     -- node[lbl, midway, above] {/cmd\_vel}          (motion.west);
\draw[flow] (battery.east) .. controls (5,0) and (7,2) .. node[lbl, pos=0.55, below] {/battery\_state} (nav.south west);
\draw[flow] (detector.east) -- node[lbl, midway, above] {/detected\_objects} (nav.north east);
\draw[flow] (state.east) -- node[lbl, midway, above] {/robot/status} (ui.south west);

% --- Legende ---
\node[ros, minimum width=8mm, minimum height=5mm, font=\small] (leg1) at (3.2,-3.4) {};
\node[right=2mm of leg1, font=\small, align=left] {ROS2-Node};
\draw[flow] (8.4,-3.4) -- ++(12mm,0);
\node[font=\small, align=left, anchor=west] at (9.8,-3.4) {Topic (Nachrichtenfluss)};

\end{tikzpicture}

}
%     \caption{Node-Graph des autonomen Indoor-Roboters mit den wichtigsten
%       Topics zwischen Sensor-, Verarbeitungs-, Navigations- und
%       Aktuatorik-Nodes.}
%     \label{fig:kap15-node-graph}
%   \end{figure}
%
% Benoetigte Pakete (im Buch bereits in main.tex geladen): tikz, graphicx
% Benoetigte TikZ-Bibliotheken: positioning, arrows.meta, shadows, calc
% =====================================================================

\begin{tikzpicture}[
    scale=0.74, transform shape,
    every node/.style={font=\small},
    node distance=10mm and 14mm,
    ros/.style={
        rectangle, rounded corners=2pt, draw=orange!60!black, thick,
        minimum width=32mm, minimum height=10mm,
        align=center, fill=orange!15, drop shadow
    },
    flow/.style={-{Stealth[length=2.2mm]}, thick, orange!70!black},
    lbl/.style={font=\footnotesize, fill=white, inner sep=1pt}
]

% --- Spalte 1: Sensor-Nodes (links) ---
\node[ros] (camera)  at (0, 6)  {\texttt{camera\_node}};
\node[ros] (lidar)   at (0, 3)  {\texttt{lidar\_node}};
\node[ros] (battery) at (0, 0)  {\texttt{battery\_monitor\_}\\\texttt{node}};

% --- Spalte 2: Verarbeitung ---
\node[ros] (detector) at (5.6, 6) {\texttt{object\_detector\_}\\\texttt{node}};
\node[ros] (slam)     at (5.6, 3) {\texttt{slam\_node}};

% --- Spalte 3: zentral ---
\node[ros] (nav) at (10.2, 3.5) {\texttt{navigation\_node}};

% --- Spalte 4: rechts ---
\node[ros] (motion) at (14.8, 5)   {\texttt{motion\_control\_}\\\texttt{node}};
\node[ros] (ui)      at (14.8, 0.5) {\texttt{user\_interface\_}\\\texttt{node}};

% --- zusätzlicher Node (Status) ---
\node[ros] (state) at (10.2, -1.5) {\texttt{robot\_state\_node}};

% --- Kanten ---
\draw[flow] (camera.east)  -- node[lbl, midway, above] {/camera/image\_raw} (detector.west);
\draw[flow] (lidar.east)   -- node[lbl, midway, above] {/scan}              (slam.west);
\draw[flow] (slam.east)    -- node[lbl, midway, above] {/map}               (nav.north west);
\draw[flow] (nav.east)     -- node[lbl, midway, above] {/cmd\_vel}          (motion.west);
\draw[flow] (battery.east) .. controls (5,0) and (7,2) .. node[lbl, pos=0.55, below] {/battery\_state} (nav.south west);
\draw[flow] (detector.east) -- node[lbl, midway, above] {/detected\_objects} (nav.north east);
\draw[flow] (state.east) -- node[lbl, midway, above] {/robot/status} (ui.south west);

% --- Legende ---
\node[ros, minimum width=8mm, minimum height=5mm, font=\small] (leg1) at (3.2,-3.4) {};
\node[right=2mm of leg1, font=\small, align=left] {ROS2-Node};
\draw[flow] (8.4,-3.4) -- ++(12mm,0);
\node[font=\small, align=left, anchor=west] at (9.8,-3.4) {Topic (Nachrichtenfluss)};

\end{tikzpicture}

}
%     \caption{Node-Graph des autonomen Indoor-Roboters mit den wichtigsten
%       Topics zwischen Sensor-, Verarbeitungs-, Navigations- und
%       Aktuatorik-Nodes.}
%     \label{fig:kap15-node-graph}
%   \end{figure}
%
% Benoetigte Pakete (im Buch bereits in main.tex geladen): tikz, graphicx
% Benoetigte TikZ-Bibliotheken: positioning, arrows.meta, shadows, calc
% =====================================================================

\begin{tikzpicture}[
    scale=0.74, transform shape,
    every node/.style={font=\small},
    node distance=10mm and 14mm,
    ros/.style={
        rectangle, rounded corners=2pt, draw=orange!60!black, thick,
        minimum width=32mm, minimum height=10mm,
        align=center, fill=orange!15, drop shadow
    },
    flow/.style={-{Stealth[length=2.2mm]}, thick, orange!70!black},
    lbl/.style={font=\footnotesize, fill=white, inner sep=1pt}
]

% --- Spalte 1: Sensor-Nodes (links) ---
\node[ros] (camera)  at (0, 6)  {\texttt{camera\_node}};
\node[ros] (lidar)   at (0, 3)  {\texttt{lidar\_node}};
\node[ros] (battery) at (0, 0)  {\texttt{battery\_monitor\_}\\\texttt{node}};

% --- Spalte 2: Verarbeitung ---
\node[ros] (detector) at (5.6, 6) {\texttt{object\_detector\_}\\\texttt{node}};
\node[ros] (slam)     at (5.6, 3) {\texttt{slam\_node}};

% --- Spalte 3: zentral ---
\node[ros] (nav) at (10.2, 3.5) {\texttt{navigation\_node}};

% --- Spalte 4: rechts ---
\node[ros] (motion) at (14.8, 5)   {\texttt{motion\_control\_}\\\texttt{node}};
\node[ros] (ui)      at (14.8, 0.5) {\texttt{user\_interface\_}\\\texttt{node}};

% --- zusätzlicher Node (Status) ---
\node[ros] (state) at (10.2, -1.5) {\texttt{robot\_state\_node}};

% --- Kanten ---
\draw[flow] (camera.east)  -- node[lbl, midway, above] {/camera/image\_raw} (detector.west);
\draw[flow] (lidar.east)   -- node[lbl, midway, above] {/scan}              (slam.west);
\draw[flow] (slam.east)    -- node[lbl, midway, above] {/map}               (nav.north west);
\draw[flow] (nav.east)     -- node[lbl, midway, above] {/cmd\_vel}          (motion.west);
\draw[flow] (battery.east) .. controls (5,0) and (7,2) .. node[lbl, pos=0.55, below] {/battery\_state} (nav.south west);
\draw[flow] (detector.east) -- node[lbl, midway, above] {/detected\_objects} (nav.north east);
\draw[flow] (state.east) -- node[lbl, midway, above] {/robot/status} (ui.south west);

% --- Legende ---
\node[ros, minimum width=8mm, minimum height=5mm, font=\small] (leg1) at (3.2,-3.4) {};
\node[right=2mm of leg1, font=\small, align=left] {ROS2-Node};
\draw[flow] (8.4,-3.4) -- ++(12mm,0);
\node[font=\small, align=left, anchor=west] at (9.8,-3.4) {Topic (Nachrichtenfluss)};

\end{tikzpicture}

